\section{مقدمه}

\subsection{تاثیر دیستروفی عضلانی دوشن بر عملکرد حرکتی}

اگرچه در سال‌های اخیر پیشرفت‌های قابل توجهی در زمینه درمان‌های دارویی، ژن‌درمانی و مراقبت‌های حمایتی صورت گرفته است، اما هنوز درمان قطعی برای این بیماری وجود ندارد. به همین دلیل، پایش مستمر وضعیت عملکردی بیماران و ارزیابی روند پیشرفت بیماری از اهمیت بالایی برخوردار است. نتایج این ارزیابی‌ها علاوه بر کمک به انتخاب روش درمان، در بررسی اثربخشی مداخلات درمانی، برنامه‌ریزی جلسات توان‌بخشی و انجام مطالعات بالینی نیز نقش اساسی ایفا می‌کنند \cite{Bushby2010Diagnosis}.

توسعه سامانه‌هایی که بتوانند اطلاعات حرکتی بیماران را به صورت کمّی ثبت کرده و با حداقل وابستگی به نظر ارزیاب تحلیل کنند، در سال‌های اخیر مورد توجه بسیاری از پژوهشگران قرار گرفته است. بهره‌گیری از حسگرهای پوشیدنی و روش‌های تحلیل داده این امکان را فراهم می‌کند که بخشی از فرآیند ارزیابی عملکرد بیماران به صورت خودکار انجام شود. پژوهش حاضر نیز با همین رویکرد طراحی و اجرا شده است.


\subsection{ارزیابی عملکرد حرکتی بیماران و چالش‌های موجود}

از آنجا که دیستروفی عضلانی دوشن یک بیماری پیشرونده است، ارزیابی مستمر وضعیت عملکرد حرکتی بیماران نقش مهمی در مدیریت روند درمان و توان‌بخشی آن‌ها ایفا می‌کند. پزشکان و درمانگران با بررسی تغییرات عملکرد حرکتی در طول زمان، می‌توانند سرعت پیشرفت بیماری را ارزیابی کرده، اثربخشی مداخلات درمانی را بررسی کنند و در صورت نیاز برنامه درمانی بیمار را اصلاح نمایند. به همین دلیل، ارزیابی عملکرد حرکتی به یکی از بخش‌های جدایی‌ناپذیر مراقبت از بیماران مبتلا به دیستروفی عضلانی دوشن تبدیل شده است \cite{Bushby2010Diagnosis}.

بر خلاف بسیاری از بیماری‌ها که وضعیت بیمار را می‌توان با یک آزمایش آزمایشگاهی یا تصویربرداری مشخص ارزیابی کرد، در دیستروفی عضلانی دوشن بخش قابل توجهی از اطلاعات مورد نیاز پزشک از نحوه انجام فعالیت‌های روزمره توسط بیمار به دست می‌آید. فعالیت‌هایی نظیر راه رفتن، برخاستن از زمین، بالا رفتن از پله‌ها، حفظ تعادل و سایر حرکات عملکردی، اطلاعات ارزشمندی درباره وضعیت عضلات و توانایی حرکتی بیمار در اختیار پزشک قرار می‌دهند.

در حال حاضر، بخش عمده این ارزیابی‌ها از طریق مشاهده مستقیم بیمار توسط پزشک یا درمانگر انجام می‌شود. اگرچه این روش‌ها طی سال‌ها اعتبار بالینی خود را اثبات کرده‌اند، اما ماهیت مشاهده‌ای آن‌ها سبب می‌شود نتایج تا حدی به تجربه ارزیاب، شرایط انجام آزمون و حتی وضعیت جسمی و روحی بیمار در زمان ارزیابی وابسته باشد. علاوه بر این، انجام این ارزیابی‌ها مستلزم حضور بیمار در مرکز درمانی بوده و معمولاً در فواصل زمانی نسبتاً طولانی تکرار می‌شوند؛ در نتیجه امکان پایش مداوم روند پیشرفت بیماری محدود خواهد بود.

در سال‌های اخیر، توسعه روش‌های کمی برای تحلیل حرکت انسان توجه بسیاری از پژوهشگران را به خود جلب کرده است. هدف این روش‌ها، استخراج اطلاعات قابل اندازه‌گیری از حرکات بدن و کاهش وابستگی ارزیابی به قضاوت ذهنی افراد است. استفاده از تجهیزات اندازه‌گیری حرکت و روش‌های پردازش داده این امکان را فراهم می‌کند که ویژگی‌های سینماتیکی حرکت به‌صورت عددی استخراج شده و تغییرات عملکرد بیمار در طول زمان با دقت بیشتری بررسی شود \cite{Winter2009Biomechanics}.

با توجه به این نیاز، پژوهش‌های متعددی در سال‌های اخیر به دنبال توسعه سامانه‌هایی برای ثبت، تحلیل و ارزیابی خودکار حرکات بیماران بوده‌اند. اگرچه در بسیاری از مطالعات از سامانه‌های اپتیکی تحلیل حرکت نیز استفاده شده است، اما این سامانه‌ها معمولاً به فضای آزمایشگاهی، دوربین‌های متعدد و فرآیند کالیبراسیون پیچیده نیاز دارند. در مقابل، حسگرهای اینرسیایی
\footnote{\lr{Inertial Measurement Unit (IMU)}}
 ضمن حفظ دقت مناسب برای بسیاری از کاربردهای بالینی، هزینه کمتر، قابلیت حمل و سهولت استفاده بیشتری دارند. این ویژگی‌ها سبب شده است که استفاده از حسگرهای اینرسیایی به یکی از رویکردهای رایج در تحلیل حرکت بیماران تبدیل شود. انتخاب ابزار اندازه‌گیری مناسب و معیار بالینی قابل اعتماد، دو مؤلفه اصلی در طراحی چنین سامانه‌هایی به شمار می‌روند. در بخش‌های بعدی، ابتدا روش‌های متداول اندازه‌گیری حرکت معرفی شده و سپس معیار ارزیابی انتخاب‌شده در این پژوهش ارائه خواهد شد.

\subsection{بیان مسئله}

ارزیابی عملکرد حرکتی یکی از مهم‌ترین ابزارهای پایش روند پیشرفت بیماری در بیماران مبتلا به دیستروفی عضلانی دوشن است. نتایج این ارزیابی‌ها در تصمیم‌گیری‌های درمانی، برنامه‌ریزی جلسات توان‌بخشی، بررسی اثربخشی داروها و همچنین انجام مطالعات بالینی مورد استفاده قرار می‌گیرد. از این رو، دقت و تکرارپذیری فرآیند ارزیابی اهمیت بالایی دارد.

در روش‌های رایج، بخش عمده‌ای از ارزیابی عملکرد بیماران توسط پزشک یا درمانگر و بر اساس مشاهده مستقیم نحوه انجام فعالیت‌های مختلف صورت می‌گیرد. اگرچه این روش‌ها طی سال‌ها اعتبار بالینی خود را اثبات کرده‌اند، اما بخشی از فرآیند نمره‌دهی به تجربه ارزیاب و برداشت وی از کیفیت اجرای حرکت وابسته است. علاوه بر این، انجام ارزیابی مستلزم حضور بیمار در مرکز درمانی بوده و امکان ثبت مداوم تغییرات عملکردی بیمار در فواصل زمانی کوتاه وجود ندارد.

از سوی دیگر، پیشرفت فناوری حسگرهای پوشیدنی و روش‌های تحلیل داده، امکان اندازه‌گیری کمی بسیاری از ویژگی‌های حرکتی را فراهم کرده است. ثبت شتاب خطی و سرعت زاویه‌ای اندام‌های مختلف بدن می‌تواند اطلاعات ارزشمندی درباره نحوه اجرای حرکات در اختیار قرار دهد. با این حال، تبدیل این داده‌های خام به شاخص‌هایی که بتوانند وضعیت عملکردی بیمار را توصیف کنند، همچنان یکی از چالش‌های اصلی این حوزه محسوب می‌شود.

بر همین اساس، مسئله اصلی این پژوهش طراحی و توسعه سامانه‌ای مبتنی بر حسگرهای اینرسیایی است که بتواند داده‌های حرکتی بیماران دیستروفی عضلانی دوشن را ثبت کرده، پس از پردازش مناسب، ویژگی‌های مؤثر را استخراج نموده و با استفاده از روش‌های یادگیری ماشین، ارتباط میان داده‌های حرکتی و ارزیابی بالینی بیماران را مدل‌سازی کند.

\subsection{اهداف پژوهش}

هدف این پژوهش توسعه یک چارچوب یکپارچه برای ثبت داده‌های حرکتی بیماران مبتلا به دیستروفی عضلانی دوشن، استخراج ویژگی‌های مناسب از سیگنال‌های ثبت‌شده و پیش‌بینی نمرات عملکردی با استفاده از الگوریتم‌های یادگیری ماشین است. انتظار می‌رود نتایج این پژوهش بتواند گامی در جهت ارائه روشی کم‌هزینه، کمّی و قابل تکرار برای کمک به ارزیابی عملکرد حرکتی بیماران باشد.

اهداف فرعی پژوهش عبارت‌اند از:

\begin{itemize}
	\item مطالعه روش‌های متداول ارزیابی عملکرد حرکتی بیماران مبتلا به دیستروفی عضلانی دوشن ;
	
	\item طراحی سامانه اندازه‌گیری مبتنی بر حسگرهای اینرسیایی و انتخاب محل مناسب نصب حسگرها؛
	
	\item توسعه نرم‌افزار دریافت، ذخیره‌سازی و مدیریت داده‌های حسگرها؛
	
	\item پیش‌پردازش داده‌های ثبت‌شده و استخراج ویژگی‌های مناسب از سیگنال‌های حرکتی؛
	
	\item بررسی تأثیر روش‌های مختلف کاهش بعد بر داده‌های استخراج‌شده؛
	
	\item آموزش و ارزیابی مدل‌های مختلف یادگیری ماشین برای پیش‌بینی نمرات عملکردی؛
	
	\item مقایسه عملکرد مدل‌های مختلف و انتخاب مناسب‌ترین مدل برای مسئله مورد مطالعه.
\end{itemize}


\subsection{خروجی‌های مورد انتظار}

انتظار می‌رود سامانه توسعه‌یافته بتواند به‌عنوان ابزاری مکمل در کنار ارزیابی‌های بالینی مورد استفاده قرار گیرد و با ارائه شاخص‌های کمی، به افزایش عینیت، دقت و تکرارپذیری فرآیند نمره‌دهی کمک نماید. همچنین این چارچوب می‌تواند زمینه‌ای برای توسعه سامانه‌های پایش طولی بیماران و کاربردهای ارزیابی از راه دور در آینده فراهم سازد.

% \subsection{ساختار گزارش}

% این گزارش در چهار فصل تنظیم شده است.

% در فصل اول، ابتدا بیماری دیستروفی عضلانی دوشن و ضرورت ارزیابی عملکرد حرکتی بیماران معرفی شده و سپس روند انتخاب روش ارزیابی، انتخاب حسگرهای اینرسیایی و بیان مسئله پژوهش ارائه شده است.

% فصل دوم به تشریح مراحل انجام پروژه اختصاص دارد. در این فصل، طراحی سامانه اندازه‌گیری، فرآیند ثبت داده‌ها، توسعه نرم‌افزار جمع‌آوری اطلاعات، پیش‌پردازش داده‌ها، استخراج ویژگی‌ها و طراحی مدل‌های یادگیری ماشین به تفصیل بیان می‌شود.

% در فصل سوم، نتایج حاصل از اجرای سامانه، عملکرد مدل‌های مختلف و تحلیل نتایج ارائه و مورد بررسی قرار می‌گیرد.

% در نهایت، فصل چهارم به جمع‌بندی پژوهش، بیان محدودیت‌ها و پیشنهادهایی برای ادامه این مسیر تحقیقاتی اختصاص یافته است.